\chapter{Kommentierte Quellen}

Achtundzwanzig Werke, eingeordnet: Zu jedem steht, wofür man es aufschlägt und was es leistet, das andere nicht leisten.

\lit{1}{Abarbanel, H.: \textit{Analysis of Observed Chaotic Data}. Springer 1996.}{Teil~IV aus der Hand eines Physikers, der die Methoden mitentwickelt hat: Einbettung, falsche Nachbarn, Exponenten aus Daten, Vorhersage. \textbf{Die Quelle für Satz 10.3 und das Benettin-Verfahren an Daten}; stellenweise optimistischer über die Diagnostik als Kapitel~12.}

\lit{2}{Argyris, J., Faust, G., Haase, M., Friedrich, R.: \textit{Die Erforschung des Chaos}. 2.\ Aufl., Springer 2010.}{\textbf{Auf Deutsch, achthundert Seiten}: Fixpunkte, Zyklen, Verzweigungen, Lorenz, Dimensionen, Einbettung, mit vielen Abbildungen und Beispielen aus der Technik. Das umfassendste deutschsprachige Werk zum Band; die Kapitel~2 bis 11 finden sich dort in größerer Breite.}

\lit{3}{Barkley, D.: \textit{Linear Stability Analysis of Rotating Spiral Waves in Excitable Media}. Phys.\ Rev.\ Lett.\ 68, 1992.}{Die Spirale als Verzweigungsgegenstand: Ihre Spitzenbewegung geht durch eine Hopf-Verzweigung vom Kreis zum Mäander. \textbf{Kapitel~5 auf Kapitel~13 angewandt}; mit dem Modell, das seither die Spiralforschung trägt.}

\lit{4}{Devaney, R.: \textit{An Introduction to Chaotic Dynamical Systems}. 2.\ Aufl., Addison-Wesley 1989.}{Die mathematische Theorie der Abbildungen --- Scharkowski mit Beweis, Schwarzsche Ableitung, symbolische Dynamik, Hufeisen. \textbf{Die Fundstelle für die Auslassungen~7 und~9}; für Leser, die Kapitel~7 mit Beweisen wollen.}

\lit{5}{Eckmann, J.-P., Ruelle, D.: \textit{Ergodic Theory of Chaos and Strange Attractors}. Rev.\ Mod.\ Phys.\ 57, 1985.}{Die Übersicht, die Exponenten, Dimensionen und Entropie in einen Rahmen stellt: Oseledets, Pesin, Kaplan--Yorke, Grassberger--Procaccia --- Kapitel~9 und~11 in der Sprache der Ergodentheorie. \textbf{Die Quelle für die Auslassungen~10, 11 und~16}; anspruchsvoll, aber die Einleitung ist für jeden lesbar.}

\lit{6}{Feigenbaum, M.: \textit{Quantitative Universality for a Class of Nonlinear Transformations}. J.\ Stat.\ Phys.\ 19, 1978.}{Satz 7.3 im Original: die Beobachtung am Rechner, die Renormierungsgleichung, die Zahl. \textbf{Lesenswert für die Entdeckung}, die vor dem Beweis stand; Feigenbaums Aufsatz in Los Alamos Science 1980 ist die zugänglichere Darstellung.}

\lit{7}{Fenton, F., Cherry, E., Hastings, H., Evans, S.: \textit{Multiple Mechanisms of Spiral Wave Breakup in a Model of Cardiac Electrical Activity}. Chaos 12, 2002.}{Der Katalog der Spiralbrüche: steile Restitution, diskordanter Alternans, Mäander, Anisotropie --- mit Simulationen für jeden. \textbf{Die Quelle für den Bruchfall in Kapitel~13} und für die Frage, warum FitzHugh--Nagumo nicht bricht.}

\lit{8}{Glass, L., Mackey, M.: \textit{From Clocks to Chaos: The Rhythms of Life}. Princeton 1988.}{Das Buch, das aus Herzzellen, Atmung und Blutbildung ein Fach machte: Phasenantwort, Kreisabbildung, Blockbilder, Periodenverdopplung. \textbf{Die Vorlage für Kapitel~6 und~7} und für den Leitfall des Bandes; ohne Beweise, mit Daten.}

\lit{9}{Grassberger, P., Procaccia, I.: \textit{Measuring the Strangeness of Strange Attractors}. Physica D 9, 1983.}{Satz 11.3 im Original: der Korrelationsalgorithmus, die Ungleichung $D_2 \le D_0$, die Werte für Lorenz und Hénon. \textbf{Zwölf Seiten, die zehntausendmal zitiert wurden} --- und deren Warnungen vor zu kurzen Reihen man überlesen hat.}

\lit{10}{Guckenheimer, J., Holmes, P.: \textit{Nonlinear Oscillations, Dynamical Systems, and Bifurcations of Vector Fields}. Springer 1983.}{Das Standardwerk der mathematischen Seite: Zentrumsmannigfaltigkeiten, Normalformen, Hufeisen, Melnikov. \textbf{Die Fundstelle für Auslassung~17} und für alles aus Teil~II, was hier nur in der Ebene gezeigt wurde. Schwer, vollständig, unersetzt.}

\lit{11}{Izhikevich, E.: \textit{Dynamical Systems in Neuroscience: The Geometry of Excitability and Bursting}. MIT Press 2007.}{Teil~I und~II für erregbare Zellen: Nullklinen, Schwellen, Hopf und Sattel-Knoten auf dem Zyklus, Phasenantwortkurven, mit Hunderten Abbildungen. \textbf{Die Zelle dieses Bandes, von einem Neurowissenschaftler erklärt}; die Klassifikation der Erregbarkeit nach dem Verzweigungstyp ist das Herzstück.}

\lit{12}{Kantz, H., Schreiber, T.: \textit{Nonlinear Time Series Analysis}. 2.\ Aufl., Cambridge 2004.}{\textbf{Das Lehrbuch zu Teil~IV}: Einbettung, Vorhersage, Exponenten, Dimensionen, Surrogate --- mit allen Fallen und dem Programmpaket TISEAN, mit dem die Methoden seither gerechnet werden. Die Quelle für die Bruchfälle der Kapitel~10 bis 12; von zwei Autoren aus Dresden.}

\lit{13}{Keener, J., Sneyd, J.: \textit{Mathematical Physiology}. 2.\ Aufl., Springer 2009.}{Zwei Bände, in denen die Mathematik dieses Bandes auf Zelle, Herz, Nerv und Hormon angewandt wird: Band~I die Zelle (FitzHugh--Nagumo, Fronten, Wellen), Band~II die Organe (Reentry, Spiralen, Defibrillation). \textbf{Die Fundstelle für Auslassung~15} und für Kapitel~13 in der Physiologie.}

\lit{14}{Kuramoto, Y.: \textit{Chemical Oscillations, Waves, and Turbulence}. Springer 1984 (Dover 2003).}{Das Modell aus Kapitel~14 in der Fassung seines Erfinders, nebst Phasenreduktion --- der Theorie, die aus jedem Oszillator eine Phase macht und Kapitel~6 begründet. \textbf{Knapp, dicht, grundlegend}; die Phasenreduktion in Kapitel~2 bis 4 ist der Teil, den jeder lesen sollte.}

\lit{15}{Kuznetsov, Y.: \textit{Elements of Applied Bifurcation Theory}. 3.\ Aufl., Springer 2004.}{Das Standardwerk der Verzweigungstheorie: Sattel-Knoten, Hopf, Flip, Zentrums\-mannig\-faltig\-keit, Normal\-formen, mit den Formeln für die Koeffizienten. \textbf{Die Quelle für Satz 5.1 und den Ljapunow-Koeffizienten}, Auslassungen~4 und~5.}

\lit{16}{Leven, R., Koch, B.-P., Pompe, B.: \textit{Chaos in dissipativen Systemen}. 2.\ Aufl., Akademie Verlag 1994.}{\textbf{Auf Deutsch, aus Greifswald}: Attraktoren, Exponenten, Dimensionen, Wege ins Chaos, mit Beweisskizzen und Übungen. Die Kapitel~7 bis 11 in deutscher Sprache, knapper als Argyris und mathematischer; die Darstellung der Dimensionen ist besonders klar.}

\lit{17}{Lorenz, E.: \textit{Deterministic Nonperiodic Flow}. J.\ Atmos.\ Sci.\ 20, 1963.}{Kapitel~8 im Original: die drei Gleichungen, der Attraktor, die Lorenz-Abbildung mit dem Argument gegen stabile Zyklen, der Schluss über die Vorhersagbarkeit. \textbf{Einer der meistzitierten Aufsätze der Naturwissenschaft}, und einer der klarsten; zwanzig Seiten, die man gelesen haben sollte.}

\lit{18}{Metzler, W.: \textit{Nichtlineare Dynamik und Chaos}. Teubner 1998.}{\textbf{Auf Deutsch}: Abbildungen und Flüsse, Verzweigungen, Fraktale, mit vielen durchgerechneten Beispielen und Rechnerexperimenten. Die zugänglichste deutschsprachige Einführung zu den Teilen~II und~III; für Leser, die mit den Übungsaufgaben dieses Bandes gut zurechtkommen.}

\lit{19}{Ott, E.: \textit{Chaos in Dynamical Systems}. 2.\ Aufl., Cambridge 2002.}{Die Physik des Chaos auf Lehrbuchniveau: Attraktoren, Dimensionen (Hausdorff, $D_q$, Multifraktale), Exponenten, Kreisabbildung und Zungen, Hamiltonsche Systeme. \textbf{Die Quelle für die Auslassungen~6, 11, 13 und~14}; das Buch, das die Breite des Fachs am besten zeigt.}

\lit{20}{Perko, L.: \textit{Differential Equations and Dynamical Systems}. 3.\ Aufl., Springer 2001.}{Die mathematischen Grundlagen von Teil~I mit Beweisen: Existenz und Eindeutigkeit, Hartman--Grobman, stabile Mannigfaltigkeiten, Poincaré--Bendixson. \textbf{Die Fundstelle für die Auslassungen~1 bis~3}; für Leser mit Band 14.}

\lit{21}{Pikovsky, A., Rosenblum, M., Kurths, J.: \textit{Synchronization: A Universal Concept in Nonlinear Sciences}. Cambridge 2001.}{Kapitel~6 und~14 als eigenes Fach: Phasenantwort, Zungen, Kuramoto, Synchronisation chaotischer Systeme, mit Beispielen aus Physik und Physiologie. \textbf{Das Standardwerk der Synchronisation}; von einem Potsdamer Autorenteam.}

\lit{22}{Schuster, H.\ G., Just, W.: \textit{Deterministic Chaos: An Introduction}. 4.\ Aufl., Wiley-VCH 2005 (deutsch: \textit{Deterministisches Chaos}, VCH 1994).}{Die Physikervariante von Kapitel~7 bis 11: Feigenbaum mit Renormierung, Intermittenz, seltsame Attraktoren, Exponenten. \textbf{Die beste Darstellung der Renormierung für Nichtmathematiker}, Auslassung~8; die deutsche Ausgabe der dritten Auflage ist noch zu finden.}

\lit{23}{Strogatz, S.: \textit{Nonlinear Dynamics and Chaos}. 2.\ Aufl., Westview 2015.}{\textbf{Das Lehrbuch, aus dem die meisten das Fach gelernt haben}: Fixpunkte, Verzweigungen, Zyklen, Lorenz, Abbildungen, Fraktale --- in der Reihenfolge dieses Bandes, mit Beispielen aus Biologie und Technik und Aufgaben, die sich lohnen. Für Teil~I bis III die erste Adresse; Teil~IV und~V fehlen.}

\lit{24}{Takens, F.: \textit{Detecting Strange Attractors in Turbulence}. In: Dynamical Systems and Turbulence, Warwick 1980. Lecture Notes in Math.\ 898, Springer 1981.}{Satz 10.2 im Original, auf vierzehn Seiten, mit dem Beweis über Whitney und Transversalität. \textbf{Die Fundstelle für Auslassung~12}; die Einleitung sagt in einem Absatz, warum ein Kanal genügt.}

\lit{25}{Theiler, J., Eubank, S., Longtin, A., Galdrikian, B., Farmer, J.\ D.: \textit{Testing for Nonlinearity in Time Series: The Method of Surrogate Data}. Physica D 58, 1992.}{Kapitel~12 im Original: Nullhypothese, FT- und AAFT-Surrogate, Rangstatistik, die Beispiele. \textbf{Der Aufsatz, der aus der Chaos-Diagnostik Statistik machte}; mit Schreiber--Schmitz (Physica D 142, 2000) als Fortsetzung.}

\lit{26}{Tucker, W.: \textit{A Rigorous ODE Solver and Smale's 14th Problem}. Found.\ Comput.\ Math.\ 2, 2002.}{Der Beweis, dass der Lorenz-Attraktor existiert --- Intervallarithmetik plus Normalform, Kapitel~8. \textbf{Nicht zum Lesen, zum Wissen, dass es ihn gibt}; Vianas Übersicht in Math.\ Intelligencer 22, 2000, erklärt die Idee.}

\lit{27}{Wiggins, S.: \textit{Introduction to Applied Nonlinear Dynamical Systems and Chaos}. 2.\ Aufl., Springer 2003.}{Das zweite große mathematische Lehrbuch neben Guckenheimer--Holmes, mit ausführlicheren Beweisen und mehr Beispielen: Normalformen, Zentrumsmannigfaltigkeiten, Hufeisen, Melnikov. \textbf{Die Fundstelle für Auslassung~4} und für Leser, die Guckenheimer--Holmes zu knapp finden.}

\lit{28}{Winfree, A.: \textit{The Geometry of Biological Time}. 2.\ Aufl., Springer 2001.}{Das Buch über Phasen, Phasenantwortkurven, Phasensingularitäten und Spiralen --- Kapitel~6 und~13 aus der Hand dessen, der beides für die Biologie erfand. \textbf{Eigenwillig, weitschweifig, unersetzlich}; die Kapitel über die circadiane Uhr und das Herz sind die Quelle der Beispiele in Aufgabe 6.3 und Kapitel~13.}
